\begin{trapbox}
\begin{description}[style=nextline, leftmargin=2.8em, labelsep=0.5em, itemsep=0.35em]
    \item[\textbf{\textcolor{CTrap}{易错点一（忽略前提条件）}}] 忘记 $x \geq -1$ 的前提，直接对任意 $x$ 使用不等式。
    \item[\textbf{\textcolor{CTrap}{正确理解（保障定义域）：}}] 条件 $x \geq -1$ 确保了底数 $1+x \geq 0$，从而 $(1+x)^n$ 对任意实数 $n$ 有意义（实数幂要求底数非负）。
    \item[\textbf{\textcolor{CTrap}{错因分析（盲目套用）：}}] 学生只记住了不等式形式，忽略了该不等式成立的前提条件，导致在 $x < -1$ 时错误使用。
\end{description}

\medskip
\begin{description}[style=nextline, leftmargin=2.8em, labelsep=0.5em, itemsep=0.35em]
    \item[\textbf{\textcolor{CTrap}{易错点二（方向记反）}}] 混淆 $n$ 在不同范围内不等号的方向，错误地认为 $n$ 越大 $(1+x)^n$ 越大，所以总是 $\geq$。
    \item[\textbf{\textcolor{CTrap}{正确理解（指数分界）：}}] 不等号方向以 $n=0$ 和 $n=1$ 为界。当 $n$ “跑出”区间 $[0,1]$ 时为 $\geq$；“跑进”区间 $[0,1]$ 时为 $\leq$。
    \item[\textbf{\textcolor{CTrap}{错因分析（直观错觉）：}}] 学生误将“指数 $n$ 越大，函数值越大”这一单调性（当 $1+x>1$ 时成立）错误推广到与线性函数的比较上，未理解不等式方向由函数的凹凸性决定。
\end{description}

\medskip
\begin{description}[style=nextline, leftmargin=2.8em, labelsep=0.5em, itemsep=0.35em]
    \item[\textbf{\textcolor{CTrap}{易错点三（形式不匹配）}}] 题目中出现 $(1+ax)^b$ 或类似形式，未将其变形为标准形式 $(1+x)^n$ 就盲目套用。
    \item[\textbf{\textcolor{CTrap}{正确理解（变形适配）：}}] 使用前必须将式子化为 $(1+\text{新变量})^\text{新指数}$ 的形式，并确保新变量 $\geq -1$。例如，比较 $(1+2x)^3$ 与 $1+6x$，可令 $t=2x$，则需 $t \geq -1$，即 $x \geq -0.5$。
    \item[\textbf{\textcolor{CTrap}{错因分析（机械套用）：}}] 学生看到形似但系数不同的式子，没有通过变量代换将其转化为标准形式，导致结论错误或遗漏变量的限制条件。
\end{description}
\end{trapbox}
